% Evans 8.1: printed pp.427--437 / English PDF pp.444--454.
% Chinese: printed pp.303--310 / scan pp.321--328.
\section{Introduction}
\label{sec:8-1}
\Needspace{45mm}
\subsection{First variation and the Euler--Lagrange equation}
For a smooth Lagrangian $L(p,z,x)$ define
\begin{equation}\label{eq:8-energy}
 I[u]=\int_U L(Du,u,x)\,dx,\qquad u-g\in W^{1,q}_0(U).
\end{equation}
Here $g\in W^{1,q}(U)$ is a lift of the boundary data. An energy is a
functional on functions; its derivative is initially a linear functional
on variations, not automatically a pointwise differential operator.
For a smooth minimizer $u$ and $\phi\in C_c^\infty(U)$, the curve
$i(s)=I[u+s\phi]$ has $i'(0)=0$. Differentiation and integration by parts give
\begin{equation}\label{eq:8-first-variation}
 \delta I[u](\phi)=\int_U L_p(Du,u,x)\cdot D\phi+L_z(Du,u,x)\phi=0,
\end{equation}
so the Euler--Lagrange equation is
\begin{equation}\label{eq:8-EL}
 -\operatorname{div}L_p(Du,u,x)+L_z(Du,u,x)=0.
\end{equation}
For smooth functions and compactly supported variations, the derivatives
are bounded on the compact ranges involved; this justifies differentiating
inside the integral. The Sobolev justification appears in \S\ref{sec:8-2}.

\Needspace{55mm}
\PDEAnnotation{偏导数与全导数}{$L_{x_i}$ 只对显式的 $x$ 变量求偏导；$\partial_i[L_{p_i}(Du,u,x)]$ 还包含 $D^2u$ 与 $Du$ 的链式项。Noether 公式中也必须保持这个区别。}
Expanding the equation shows its quasilinear principal part:
\[
 -L_{p_ip_j}u_{ij}-L_{p_i z}u_i-L_{p_i x_i}+L_z=0.
\]
Thus convexity of $L$ in the gradient variable determines ellipticity.
The standard examples, in the source's order, are
\[
 \begin{array}{c|c}
 L(p,z,x)&\text{Euler--Lagrange equation}\\\hline
 \tfrac12|p|^2&-\Delta u=0\\
 \tfrac12a^{ij}(x)p_i p_j-f(x)z&-\partial_i(a^{ij}u_j)=f\\
 \tfrac12|p|^2-F(z),\ F'=f&-\Delta u=f(u)\\
 (1+|p|^2)^{1/2}&\operatorname{div}\dfrac{Du}{\sqrt{1+|Du|^2}}=0
 \end{array}
\]
For the second example $A=A^T$; no arbitrary drift is generated by a
quadratic unweighted Dirichlet energy. The last integrand measures graph
area. Its Hessian has eigenvalues $(1+|p|^2)^{-1/2}$ transverse to $p$
and $(1+|p|^2)^{-3/2}$ parallel to $p$: it is strictly convex, but not
uniformly convex on all gradient space.
\Needspace{55mm}
\PDEAnnotation{极小曲面的退化}{面积密度在大梯度方向的最小曲率约为 $|p|^{-3}$。严格凸不能代替一致椭圆；面积能量还只有线性增长，不属于后面的超线性强制性框架。}

\Needspace{45mm}
\subsection{Second variation and the Legendre condition}
For a smooth local minimizer, $i''(0)\geq0$, hence
\begin{equation}\label{eq:8-second-variation}
 Q_u[\phi]=\int_U L_{p_ip_j}\phi_i\phi_j
       +2L_{p_i z}\phi_i\phi+L_{zz}\phi^2\,dx\geq0.
\end{equation}
All coefficients are evaluated at $(Du,u,x)$.
\Needspace{55mm}
\begin{theorem}[Necessary gradient convexity; source \S8.1.3]
For every $x\in U$ and $\xi\in\R^n$,
$L_{p_ip_j}(Du(x),u(x),x)\xi_i\xi_j\geq0$.
\end{theorem}
\begin{expandedproof}[small amplitude, oscillatory gradient]
Let $\rho$ be the periodic triangular function of period one, with
$|\rho'|=1$ a.e., and take
$\phi_\varepsilon=\varepsilon\rho(x\cdot\xi/\varepsilon)\zeta(x)$,
$\zeta\in C_c^\infty(U)$. Smooth approximation makes
\eqref{eq:8-second-variation} valid for this compactly supported Lipschitz
function. Its value is $O(\varepsilon)$ and
$D\phi_\varepsilon=\rho'\xi\zeta+O(\varepsilon)$, uniformly on the fixed
support except at the harmless corners. The mixed and zeroth-order terms
vanish as $\varepsilon\to0$, while $(\rho')^2=1$ a.e. gives
$\int L_{p_ip_j}\xi_i\xi_j\zeta^2\geq0$.
If this continuous coefficient were negative at a point, choosing $\zeta$
in a small neighborhood would contradict the inequality.
This is positivity along the actual gradient of $u$, not global convexity
of $L$ at all $(p,z,x)$, and it is necessary rather than sufficient for a minimum.
\end{expandedproof}
\Needspace{55mm}
\PDEAnnotation{小振幅仍有大梯度}{$\phi_\varepsilon=O(\varepsilon)$，但 $D\phi_\varepsilon=O(1)$。高频测试分离出二阶变分的主部；仅看函数值的接近不能忽略梯度振荡。}

\Needspace{45mm}
\subsection{Systems, null Lagrangians, and the Piola identity}
For $u:U\to\R^m$ write $P_{\alpha i}=\partial_i u^\alpha$.
The first variation uses the Frobenius product $A:B=\sum_{\alpha,i}A_{\alpha i}B_{\alpha i}$:
\[
 \delta I[u](\phi)=\int_U L_P:D\phi+L_z\cdot\phi,
 \qquad -\partial_i L_{P_{\alpha i}}+L_{z_\alpha}=0.
\]
A null Lagrangian is one whose Euler--Lagrange expression vanishes for
every smooth map. It need not have constant density or zero integral.
\Needspace{55mm}
\begin{theorem}[Boundary dependence; source \S8.1.4, Theorem 1]
If $L$ is null and $u,v\in C^2(\overline U;\R^m)$ have the same boundary
values, then $I[u]=I[v]$.
\end{theorem}
\begin{expandedproof}[a homotopy with fixed trace]
Put $u_s=v+s(u-v)$ for $0\leq s\leq1$. Its energy derivative is
$\int L_P(Du_s,u_s,x):D(u-v)+L_z(Du_s,u_s,x)\cdot(u-v)$.
Integration by parts has zero boundary term because $u-v=0$ there.
The remaining Euler--Lagrange expression vanishes by nullity, so
$\frac d{ds}I[u_s]=0$. Integrating in $s$ proves the claim.
\end{expandedproof}

\Needspace{55mm}
\begin{lemma}[Piola identity; source divergence-free rows lemma]
For smooth $u:U\to\R^n$, with the cofactor matrix in the same row-column
convention as $Du$,
\begin{equation}\label{eq:8-piola}
 \partial_i(\operatorname{cof}Du)_{\alpha i}=0.
\end{equation}
Consequently $L(P)=\det P$ is a null Lagrangian (source Theorem 2).
\end{lemma}
\begin{expandedproof}[antisymmetry rather than invertibility]
The entry $(\operatorname{cof}Du)_{\alpha i}$ is
\[
 \frac1{(n-1)!}\epsilon_{\alpha\alpha_2\cdots\alpha_n}
 \epsilon_{ii_2\cdots i_n}
 u^{\alpha_2}_{i_2}\cdots u^{\alpha_n}_{i_n}.
\]
Differentiate in $i$ and sum. In each resulting term the second derivative
is symmetric in $i,i_j$, whereas the spatial alternating symbol is
antisymmetric in those indices; paired terms cancel.
For $n=1$ the cofactor is the constant $1$.
Since $\partial(\det P)/\partial P_{\alpha i}=(\operatorname{cof}P)_{\alpha i}$,
\eqref{eq:8-piola} is precisely the Euler--Lagrange identity.
No division by $\det Du$ occurs, so singular gradients are included.
\end{expandedproof}
\Needspace{55mm}
\PDEAnnotation{余子式的转置约定}{这里 $Du$ 的行是分量、列是空间方向；Piola 恒等式对每行取散度。矩阵恒等式为 $P(\operatorname{cof}P)^T=(\det P)I$，不能混用伴随矩阵与余子式矩阵。}
In two dimensions it reads
$\operatorname{cof}Du=\left(\begin{smallmatrix}u^2_y&-u^2_x\\-u^1_y&u^1_x\end{smallmatrix}\right)$;
each row divergence is equality of mixed partial derivatives.
All minors enjoy analogous cancellation, a fact used in the weak continuity
argument below.

\Needspace{45mm}
\subsection{An analytic route to Brouwer's fixed point theorem}
\Needspace{55mm}
\begin{theorem}[Brouwer; source \S8.1.4, Theorem 3]
Every continuous map $F:\overline B\to\overline B$, $B=B(0,1)\subset\R^n$,
has a fixed point.
\end{theorem}
\begin{expandedproof}[no retraction, including the smoothing step]
A smooth retraction $r:\overline B\to\partial B$ would have the same
boundary data as the identity. Nullity gives $\int_B\det Dr=|B|$.
But $|r|^2=1$ implies $(Dr)^Tr=0$, so $\det Dr=0$, a contradiction.
If a continuous retraction existed, extend it by $r(x)=x$ outside $B$.
It is nonzero everywhere and equals the identity outside a compact set.
Convolution with an even radial mollifier gives smooth $r_\varepsilon$
that remains nonzero for sufficiently small $\varepsilon$: uniform
approximation on a fixed compact ball preserves the positive minimum
of $|r|$, and farther away convolution reproduces $x$ exactly.
For $|x|\geq2$, $r_\varepsilon(x)=x$ if $\varepsilon<1$.
Thus $2r_\varepsilon/|r_\varepsilon|$ is a smooth retraction from
$\overline{B(0,2)}$ onto its boundary, the same contradiction after rescaling.
Finally, if $F$ had no fixed point, the ray from $F(x)$ through $x$ would
meet the sphere in a continuous point $r(x)$ with $r(x)=x$ on $\partial B$.
Explicitly, for $a=F(x)$ and $d=x-a\ne0$, take
\[
 r(x)=a+\lambda d,\quad
 \lambda=\frac{-a\cdot d+\sqrt{(a\cdot d)^2+|d|^2(1-|a|^2)}}{|d|^2}.
\]
The forward exit parameter is at least one; the formula is continuous
because $d\ne0$. It supplies the forbidden continuous retraction.
\end{expandedproof}
